\subsection{Aprovechamiento académico}
Reconoce el esfuerzo y la dedicación durante toda la trayectoria escolar. (Art. 14, REP).

\subsubsection*{Requisitos}
\begin{itemize}
  \item Promedio general mayor o igual a 8.5.
  \item Haber cursado el plan de estudios sin interrupciones.
  \item Sin calificaciones reprobatorias ni anotaciones literales.
  \item Haber cubierto al menos el 50 \% de los créditos en la Universidad (si hubo revalidación o convalidación).
  \item Solicitar la opción por escrito dentro del primer año posterior a concluir la carrera.
  \item No haber sido sancionado por faltas a la responsabilidad universitaria.
\end{itemize}
\subsubsection*{Etapa 1. Tus datos}
Captura una sola vez tus datos; con ellos se llenan automáticamente todos los formatos del trámite.

Datos que se capturan:
\begin{itemize}
  \item Nombre completo (con acentos, tal como aparecerá en el título)
  \item Número de cuenta
  \item Licenciatura
  \item Sistema
  \item Plan de estudios (clave y versión)
  \item Escuela incorporada a la UAEM donde cursaste la licenciatura
  \item Año de ingreso
  \item Año de egreso
  \item Teléfono
  \item Correo institucional
  \item Correo personal
  \item Ocupación actual
  \item Fecha de nacimiento
  \item Mi nombre o apellidos llevan acentos
\end{itemize}
\subsubsection*{Etapa 2. Solicitud de la opción}
La solicitud debe presentarse dentro del primer año posterior a concluir tu plan de estudios.

La plataforma llena los formatos: aprovechamiento y los firma por ti.

Documentos (PDF):
\begin{longtable}{p{0.55\linewidth}p{0.35\linewidth}}
\toprule
\textbf{Documento} & \textbf{Observaciones}\\
\midrule
\endhead
Solicitud firmada por el pasante & La plataforma lo genera. La plataforma lo genera con tus datos; lo firmas electrónicamente.\\
Certificado de estudios de licenciatura & Descargado del Portal de alumnos.\\
Historial académico (trayectoria) & Firmado y sellado por el Jefe del Departamento de Control Escolar.\\
Certificado de servicio social & \\
\bottomrule
\end{longtable}
\subsubsection*{Etapa 3. Documentación para la toma de protesta}
Sube cada documento en PDF, legible y a color. Cuando todo esté validado, Apoyo a Titulación te citará para entregar originales y copias y fijar la fecha.

La plataforma llena los formatos: datos\_titulo, requisitos y los firma por ti.

Documentos (PDF):
\begin{longtable}{p{0.55\linewidth}p{0.35\linewidth}}
\toprule
\textbf{Documento} & \textbf{Observaciones}\\
\midrule
\endhead
Certificado de estudios de bachillerato & \\
Constancia de autenticación de bachillerato & Se tramita en la escuela donde cursaste el bachillerato. Solo si tus estudios de bachillerato no fueron en instituciones dependientes o incorporadas a la UAEM.\\
Acta de nacimiento & Copia certificada vigente (del año en curso).\\
CURP & Debe mostrar la leyenda «Validada por el Registro Civil».\\
Constancia de no adeudo a la Biblioteca de la Facultad & El recibo de pago se tramita en Control Escolar de la Facultad; la constancia, en la Biblioteca.\\
Constancia de no adeudo a la Biblioteca Central & \\
Contra-recibo de la donación de libros a la Biblioteca de la Facultad & \\
Constancia de no adeudo a la Dirección de Control Escolar de la UAEM & Se solicita en el Portal de alumnos: controlescolar.uaemex.mx.\\
Recibo de no adeudo a la UAEM & Se solicita en el Portal de alumnos de Control Escolar.\\
Formato de datos personales (para el título) & La plataforma lo genera. La plataforma lo llena con tus datos; lo firmas electrónicamente.\\
Comprobante del sistema de Seguimiento de Egresados & Se obtiene en alumniversitario.uaemex.mx/1dae.php.\\
Recibo de pago: derechos de evaluación profesional (\$230) & Se solicita en Control Escolar de la Facultad.\\
Recibo de pago: expedición de título (\$1,187) & Se solicita en Control Escolar de la Facultad.\\
\bottomrule
\end{longtable}
\subsubsection*{Etapa 4. Fecha, hora y lugar de la toma de protesta}
Apoyo a Titulación asigna el jurado y te comunica la fecha. Preséntate puntualmente.

\emph{Turno del departamento:} Fecha y lugar de la toma de protesta. Titulación fija la fecha. En esta opción la comisión es siempre el Director (presidente) y el Subdirector Académico (secretario); el acta oficial se genera en otro sistema y se firma en papel.
\subsubsection*{Etapa 5. Resultado de la evaluación}
Al concluir, el departamento registra el resultado. Después podrás consultar cómo descargar tu título y cédula profesional.

\emph{Turno del departamento:} Registrar resultado. 

\subsubsection*{Seguimiento y conclusión}
\subsubsection*{Pantallas del proceso}
\captura{aprovechamiento_pasante_01_opciones}{Opciones de titulación en «Licenciatura».}
\captura{aprovechamiento_pasante_02_requisitos}{Requisitos de la opción y aceptación para iniciar.}
\captura{aprovechamiento_pasante_03_iniciar}{Aceptación de condiciones para iniciar el trámite.}
\captura{aprovechamiento_pasante_04_datos__1}{Captura de datos del pasante y apertura del trámite. (parte 1 de 2)}
\captura{aprovechamiento_pasante_04_datos__2}{Captura de datos del pasante y apertura del trámite. (parte 2 de 2)}
\captura{aprovechamiento_pasante_05_ronda0__1}{Etapa en curso del trámite. (parte 1 de 2)}
\captura{aprovechamiento_pasante_05_ronda0__2}{Etapa en curso del trámite. (parte 2 de 2)}
\captura{aprovechamiento_pasante_06_concluido__1}{Trámite concluido. (parte 1 de 3)}
\captura{aprovechamiento_pasante_06_concluido__2}{Trámite concluido. (parte 2 de 3)}
\captura{aprovechamiento_pasante_06_concluido__3}{Trámite concluido. (parte 3 de 3)}
\captura{aprovechamiento_pasante_07_ruta}{Ruta del trámite con los pasos hechos.}
\captura{aprovechamiento_pasante_08_estado}{Estado del trámite.}
